%% ---------------------------------------------------------------------------
%% Deterministic sparsity and symbolic rigidity of individual Syracuse orbits
%% Build:  latexmk -pdf -outdir=build syracuse_sparsity.tex
%% ---------------------------------------------------------------------------
\documentclass[11pt,reqno]{amsart}

\usepackage[T1]{fontenc}
\usepackage[utf8]{inputenc}
\usepackage[english]{babel}
\usepackage{amsmath,amssymb,amsthm,mathtools}
\usepackage{newtxtext}
\usepackage[varg]{newtxmath}
\usepackage{microtype}
\usepackage{booktabs}
\usepackage{siunitx}
\usepackage{enumitem}
\usepackage[margin=1in]{geometry}
\usepackage{caption}
\usepackage[autostyle]{csquotes}
\usepackage[colorlinks=true,linkcolor=black,citecolor=black,urlcolor=blue]{hyperref}
\hypersetup{
  pdftitle={Deterministic sparsity and symbolic rigidity of individual Syracuse orbits},
  pdfauthor={Collaborative research draft; OpenAI Sol 5.6},
  pdfsubject={Collatz conjecture}}
\usepackage[backend=biber,style=numeric,sorting=nyt,giveninits=true,
            maxbibnames=6,doi=true,isbn=false,url=false,eprint=true]{biblatex}
\addbibresource{references.bib}

\sisetup{group-separator={\,},group-minimum-digits=4,detect-all}
\setlist[enumerate]{itemsep=2pt,topsep=4pt}
\setlist[itemize]{itemsep=2pt,topsep=4pt}

\theoremstyle{plain}
\newtheorem{theorem}{Theorem}[section]
\newtheorem{proposition}[theorem]{Proposition}
\newtheorem{corollary}[theorem]{Corollary}
\newtheorem{lemma}[theorem]{Lemma}
\theoremstyle{definition}
\newtheorem{definition}[theorem]{Definition}
\newtheorem{remark}[theorem]{Remark}
\numberwithin{equation}{section}

\newcommand{\N}{\mathbb{N}}
\newcommand{\Z}{\mathbb{Z}}
\newcommand{\R}{\mathbb{R}}
\newcommand{\Odd}{\mathcal{O}}
\newcommand{\Ent}{\mathsf{H}}

\begin{document}

\title[Sparsity and rigidity of individual Syracuse orbits]
      {Deterministic sparsity and symbolic rigidity\\ of individual Syracuse orbits}
\subjclass[2020]{Primary 11B83; Secondary 37P05, 68R15, 11K16}
\keywords{Collatz conjecture, Syracuse map, stopping time, factor complexity,
  Diophantine exponent, cylinder sets}
\date{20 August 2026}

\author{Collaborative research draft}
\author{OpenAI Sol 5.6}

\begin{abstract}
Let $S(n)=(3n+1)/2^{v_2(3n+1)}$ be the Syracuse map on odd positive integers.
The known unconditional results on the $3x+1$ problem---Terras's density-one
stopping-time theorem and Tao's almost-bounded-values theorem---are statements
about the set of starting values.  We prove four statements about
\emph{every individual} aperiodic orbit, with absolute constants.  First, such
an orbit meets a dyadic block $[X,2X)$ at most $11X^{1-\eta}$ times, where
$\eta=1-\log_23/c_*=0.0346155582\ldots$ and $c_*=1.6417941207\ldots$ is the
root of $c\log c-(c-1)\log(c-1)=\log3$; the exponent arises from an exact
balance between a Chernoff bound for valuation weights and the entropy of
valuation words.  Second, the reciprocal sum along any aperiodic orbit is
bounded by an absolute constant, so in particular every aperiodic orbit
diverges to infinity and its normalized version converges.  Third, the
valuation word of such an orbit has prefix repetition exponent at most
$\log_23$.  Fourth, its factor complexity satisfies
$p(m)\ge m/\log_2(3/2)-O(1)$, so it is never quasi-Sturmian.  All four
results follow from one cylinder lemma together with the arithmetic
non-collision of orbit values; no probabilistic or density hypothesis is used.
We do not prove, and do not claim, any progress on the Collatz conjecture
itself, and we state precisely what these bounds cannot do.
\end{abstract}

\maketitle

\tableofcontents

% ===========================================================================
\section{Introduction}
\label{sec:intro}

Write $\Odd$ for the set of odd positive integers and let
\[
S:\Odd\to\Odd,\qquad
S(n)=\frac{3n+1}{2^{a(n)}},\qquad a(n):=v_2(3n+1)\ge1
\]
be the Syracuse (odd-only Collatz) map.  The Collatz conjecture is equivalent
to the assertion that every $S$-orbit reaches $1$.

Two unconditional theorems dominate the modern literature.  Terras
\cite{terras1976} proved that the set of $n$ with finite stopping time has
density one, and Tao \cite{tao2022} proved that for every function
$f\to\infty$, almost every $n$ in logarithmic density has
$\min_k C^k(n)\le f(n)$.  Computation has verified convergence below $2^{71}$
\cite{barina2025}.  Each of these is a statement about the \emph{set of
starting values}: a density-one or almost-all conclusion says nothing about
any particular orbit, and a single exceptional orbit would suffice to refute
the conjecture.

This note isolates what can be said about one orbit at a time.  Call an orbit
$(n_k)_{k\ge0}$, $n_{k+1}=S(n_k)$, \emph{aperiodic} if $n_i\ne n_j$ whenever
$i\ne j$; since $S$ is a map, this is the same as saying the orbit is not
eventually periodic, and it is exactly the situation of a hypothetical
divergent counterexample.  Put
\[
V_s(X)=\#\{k\ge s: X\le n_k<2X\}.
\]

\begin{theorem}[shell sparsity]
\label{thm:A}
Let $c_*=1.6417941207\ldots$ be the unique root in $(\log_23,2)$ of
\[
c\log c-(c-1)\log(c-1)=\log3,
\]
and set
\[
\eta=1-\frac{\log_23}{c_*}=0.0346155582\ldots.
\]
Then for every aperiodic orbit, every $s\ge0$ and every $X\ge2$,
\[
V_s(X)\le 11\,X^{1-\eta}.
\]
The constant is absolute: it does not depend on the orbit, on $n_0$, or on $s$.
\end{theorem}

\begin{corollary}[reciprocal summability]
\label{cor:B}
For every aperiodic orbit,
\[
\sum_{k\ge0}\frac1{n_k}\le\frac{11}{1-2^{-\eta}}<464.
\]
Consequently $n_k\to\infty$, the normalized orbit
$Y_k=2^{A_k}n_k/3^k$ increases to a finite limit
$L=n_0\prod_{j\ge0}(1+\tfrac1{3n_j})<\infty$, and $n_k\le L\,2^{\log_23\,k-A_k}$
for all $k$.  If moreover $n_s=\min_{k\ge s}n_k$, then
\[
\sum_{k\ge s}\frac1{n_k}\le\frac{11}{1-2^{-\eta}}\;n_s^{-\eta}.
\]
\end{corollary}

The valuation word $a=(a_0,a_1,\dots)$, $a_k=a(n_k)$, is correspondingly
rigid.  Recall that the prefix repetition exponent (or Diophantine exponent)
$\operatorname{Dio}(a)$ is the supremum of the ratios $|P|/(|U|+|V|)$ over
prefixes $P=UV^tV'$ with $V'$ a prefix of $V$, and that $p_a(m)$ denotes the
number of distinct factors of length $m$.

\begin{theorem}[repetition and complexity]
\label{thm:CD}
Let $(n_k)$ be aperiodic with limit constant $L$ as in
Corollary~\textup{\ref{cor:B}}.  Then
\[
\operatorname{Dio}(a)\le\log_23=1.5849625\ldots,
\]
and for every $m\ge1$,
\[
p_a(m)\ \ge\ \frac{m-\log_2L}{\log_2(3/2)}-1 .
\]
In particular $\liminf_{m\to\infty}p_a(m)/m\ge1/\log_2(3/2)=1.7095112\ldots$,
so the valuation word of an aperiodic orbit is never Sturmian and never
quasi-Sturmian.
\end{theorem}

The mechanism behind all four statements is the same and is entirely
elementary.  A valuation word of length $r$ and weight $B$ pins its starting
value to a single residue class modulo $2^{B+1}$ (Lemma~\ref{lem:cyl}).
Aperiodicity says orbit values never collide.  Confronting these two facts
with the growth law of the orbit converts every hypothetical structure---too
many visits to a shell, a long repetition, a repeated factor---into an
arithmetic collision.  We use no probabilistic model of the valuations: the
geometric distribution below is a bookkeeping device for counting residue
classes exactly, not an assumption about the orbit.

Section~\ref{sec:prelim} fixes notation and proves the cylinder lemma.
Section~\ref{sec:sparsity} proves Theorem~\ref{thm:A} and
Section~\ref{sec:recip} deduces Corollary~\ref{cor:B}.
Section~\ref{sec:symbolic} proves Theorem~\ref{thm:CD}.
Section~\ref{sec:limits} states what these bounds do not give; we regard that
section as the honest counterweight to the rest of the paper.
Appendix~\ref{app:numeric} records the numerical verification of every lemma.

\begin{remark}
Nothing here bears on the existence of nontrivial \emph{cycles}, which are
periodic and therefore excluded by hypothesis throughout.
\end{remark}

% ===========================================================================
\section{Notation and the cylinder lemma}
\label{sec:prelim}

Throughout, $(n_k)_{k\ge0}$ is an orbit of $S$ in $\Odd$, and
\[
a_k=a(n_k),\qquad A_k=\sum_{j<k}a_j,\qquad A_0=0,\qquad
\alpha=\log_23 .
\]
For a finite word $w=(w_0,\dots,w_{r-1})\in\Z_{\ge1}^r$ we write
$|w|=r$ for its length and $A(w)=\sum_iw_i$ for its \emph{weight}.  We say
$n\in\Odd$ \emph{realizes} $w$ if the first $r$ valuations of the orbit of
$n$ are $w_0,\dots,w_{r-1}$.  All logarithms without a subscript are natural.

\begin{lemma}[unrolled orbit identity]
\label{lem:unroll}
For every $k\ge0$,
\begin{equation}\label{eq:unroll}
2^{A_k}n_k=3^kn_0+\sum_{j=0}^{k-1}3^{k-1-j}2^{A_j}.
\end{equation}
\end{lemma}

\begin{proof}
Induction on $k$.  The case $k=0$ is trivial.  Multiplying
$2^{a_k}n_{k+1}=3n_k+1$ by $2^{A_k}$ and substituting the inductive
hypothesis gives
$2^{A_{k+1}}n_{k+1}=3\cdot2^{A_k}n_k+2^{A_k}$, which is the claim for $k+1$.
\end{proof}

\begin{lemma}[cylinder lemma]
\label{lem:cyl}
Let $w\in\Z_{\ge1}^r$ have weight $B=A(w)$.  The set of $n\in\Odd$ realizing
$w$ is exactly one residue class modulo $2^{B+1}$.
\end{lemma}

\begin{proof}
Induction on $r$.  For $r=1$ and $w=(a)$: since $n$ is odd and $3$ is
invertible modulo $2^{a+1}$,
\[
v_2(3n+1)=a
\iff 3n+1\equiv2^a\!\!\pmod{2^{a+1}}
\iff n\equiv3^{-1}(2^a-1)\!\!\pmod{2^{a+1}},
\]
one class modulo $2^{a+1}=2^{B+1}$, and it consists of odd numbers because
$2^a-1$ is odd.

For $r\ge2$ write $w=(a_0,w')$ with $A(w')=B'=B-a_0$.  By the case $r=1$ the
first letter confines $n$ to a class $n\equiv\rho\pmod{2^{a_0+1}}$; write
$n=\rho+2^{a_0+1}t$ with $t\in\Z$.  With $m=(3\rho+1)/2^{a_0}$ odd,
\[
3n+1=3\rho+1+3\cdot2^{a_0+1}t=2^{a_0}\bigl(m+6t\bigr),
\qquad\text{so}\qquad
S(n)=m+6t .
\]
By the inductive hypothesis $S(n)$ realizes $w'$ if and only if
$S(n)\equiv\sigma\pmod{2^{B'+1}}$ for a single odd $\sigma$.  Since $m$ and
$\sigma$ are both odd, $\sigma-m=2u$ is even, and
\[
m+6t\equiv\sigma\!\!\pmod{2^{B'+1}}
\iff 3t\equiv u\!\!\pmod{2^{B'}}
\iff t\equiv3^{-1}u\!\!\pmod{2^{B'}} ,
\]
one class of $t$ modulo $2^{B'}$, hence one class of $n$ modulo
$2^{a_0+1+B'}=2^{B+1}$.
\end{proof}

\begin{lemma}[non-collision]
\label{lem:collide}
Let $(n_k)$ be aperiodic and suppose
$(a_i,\dots,a_{i+m-1})=(a_j,\dots,a_{j+m-1})$ for some $i\ne j$ and $m\ge1$.
Let $W=A_{i+m}-A_i$ be the common weight of that factor.  Then
\[
|n_i-n_j|\ge2^{W+1}.
\]
\end{lemma}

\begin{proof}
Both $n_i$ and $n_j$ realize the same word of weight $W$, so by
Lemma~\ref{lem:cyl} they lie in one class modulo $2^{W+1}$, i.e.\
$2^{W+1}\mid n_i-n_j$.  Aperiodicity gives $n_i\ne n_j$.
\end{proof}

Finally we record the two elementary estimates used in the counting argument.
Let $\Ent(x)=-x\log x-(1-x)\log(1-x)$ be the natural entropy and put
\[
g(c)=c\,\Ent(1/c)=c\log c-(c-1)\log(c-1),
\qquad
I(c)=c\log2-g(c) .
\]

\begin{lemma}[Chernoff bound for valuation weight]
\label{lem:chernoff}
Let $G_1,\dots,G_r$ be i.i.d.\ with $\Pr(G=a)=2^{-a}$, $a\ge1$.  For
$1<c<2$,
\[
\Pr\bigl(G_1+\cdots+G_r\le cr\bigr)\le e^{-rI(c)} .
\]
\end{lemma}

\begin{proof}
For $\lambda\ge0$ put $t=e^{-\lambda}\in(0,1]$.  Then
$\mathbb E[e^{-\lambda G}]=t/(2-t)$, so by Markov's inequality
\[
\Pr\Bigl(\sum G_i\le cr\Bigr)
\le e^{\lambda cr}\Bigl(\frac t{2-t}\Bigr)^r
=\Bigl(\frac{t^{1-c}}{2-t}\Bigr)^{r}.
\]
The bracket is minimized at $t=2(c-1)/c$, which lies in $(0,1)$ precisely
because $1<c<2$; there $2-t=2/c$ and the bracket equals
$(c/2)\,t^{1-c}$.  Taking logarithms,
\[
-\log\Bigl(\frac c2t^{1-c}\Bigr)
=c\log2-c\log c+(c-1)\log(c-1)=I(c). \qedhere
\]
\end{proof}

\begin{lemma}[word count]
\label{lem:words}
For integers $1\le r\le M$, $\ \sum_{B=r}^{M}\binom{B-1}{r-1}=\binom Mr$ and
$\binom Mr\le e^{M\Ent(r/M)}$.  Moreover $m\mapsto m\Ent(r/m)$ is
nondecreasing on $[r,\infty)$, so for $M=\lfloor cr\rfloor$ one has
$\binom{M}{r}\le e^{r\,g(c)}$.
\end{lemma}

\begin{proof}
The identity is the hockey-stick identity, and
$\binom Mr\le e^{M\Ent(r/M)}$ is the standard entropy bound.  For the
monotonicity, $\frac{d}{dm}\bigl(m\Ent(r/m)\bigr)=\log\frac m{m-r}>0$.
Since $\lfloor cr\rfloor\le cr$, we get
$\binom Mr\le e^{cr\Ent(1/c)}=e^{rg(c)}$.
\end{proof}

Note that the geometric law in Lemma~\ref{lem:chernoff} is not a hypothesis
about the orbit.  It enters only through the identity
$\Pr(G_1+\cdots+G_r=B)=\binom{B-1}{r-1}2^{-B}$, which is precisely the
weighted count of residue classes supplied by Lemma~\ref{lem:cyl}.

% ===========================================================================
\section{Shell sparsity}
\label{sec:sparsity}

\begin{proof}[Proof of Theorem~\textup{\ref{thm:A}}]
Since $V_s(X)\le V_0(X)$ we may assume $s=0$ and write $V=V_0$.  Let
$c=c_*$ and
\[
\kappa=\frac1{c\log2}=0.8787307876\ldots,
\qquad
r=\lceil\kappa\log X\rceil\ge1 .
\]
Two consequences of this choice are used repeatedly.  First
$c\kappa\log2=1$, so
\begin{equation}\label{eq:rlower}
2^{cr}\ge2^{c\kappa\log X}=e^{\log X}=X .
\end{equation}
Second, $r\le\kappa\log X+1$, and since $g(c_*)=\log3$ by definition of
$c_*$,
\begin{equation}\label{eq:rupper}
3^r\le3\cdot e^{\kappa\log3\,\log X}=3\,X^{\kappa\log3}=3\,X^{\alpha/c_*}
=3\,X^{1-\eta},
\end{equation}
where the last two equalities use $\kappa\log 3=\log_23/c_*=\alpha/c_*$ and
the definition of $\eta$.

Let $N=\{k\ge0:X\le n_k<2X\}$ and split $N=N_{\mathrm{low}}\cup
N_{\mathrm{high}}$ according as the block weight
$B_k:=A_{k+r}-A_k$ satisfies $B_k\le cr$ or $B_k>cr$.

\smallskip
\emph{The low-weight part.}  If $k\in N_{\mathrm{low}}$ then $n_k$ realizes a
word of length $r$ and weight $B\in[r,\lfloor cr\rfloor]$, and by
Lemma~\ref{lem:cyl} each such word confines $n_k$ to one class modulo
$2^{B+1}$, which meets $[X,2X)$ in at most $X2^{-(B+1)}+1$ integers.  By
aperiodicity distinct indices carry distinct values, so summing over words,
\[
\#N_{\mathrm{low}}
\le\sum_{B=r}^{\lfloor cr\rfloor}\binom{B-1}{r-1}
\Bigl(\frac X{2^{B+1}}+1\Bigr)
=\frac X2\sum_{B\le cr}\binom{B-1}{r-1}2^{-B}
+\sum_{B=r}^{\lfloor cr\rfloor}\binom{B-1}{r-1}.
\]
The first sum is $\Pr(G_1+\cdots+G_r\le cr)$, bounded by $e^{-rI(c_*)}$
in Lemma~\ref{lem:chernoff}; since $g(c_*)=\log3$ we have
$I(c_*)=c_*\log2-\log3$, and therefore, using $r\ge\kappa\log X$,
\[
\frac X2e^{-rI(c_*)}
\le\frac X2X^{-\kappa(c_*\log2-\log3)}
=\frac X2X^{-(1-\kappa\log3)}
=\tfrac12X^{1-\eta},
\]
because $\kappa(c_*\log2-\log3)=1-\alpha/c_*=\eta$.  The second sum equals
$\binom{\lfloor c r\rfloor}r\le e^{rg(c_*)}=3^r\le3X^{1-\eta}$ by
Lemma~\ref{lem:words} and \eqref{eq:rupper}.  Hence
\[
\#N_{\mathrm{low}}\le\tfrac72X^{1-\eta}.
\]

\smallskip
\emph{The high-weight part.}  Let $k\in N_{\mathrm{high}}$ and apply
Lemma~\ref{lem:unroll} to the block of length $r$ starting at $k$, writing
$A_{k,j}=A_{k+j}-A_k$ so that $A_{k,r}=B_k$:
\[
n_{k+r}=\frac{3^rn_k}{2^{B_k}}
+\sum_{j=0}^{r-1}3^{r-1-j}2^{-(B_k-A_{k,j})} .
\]
Every letter is at least $1$, so $B_k-A_{k,j}\ge r-j$, and substituting
$m=r-j$,
\[
\sum_{j=0}^{r-1}3^{r-1-j}2^{-(B_k-A_{k,j})}
\le\sum_{m=1}^{r}3^{m-1}2^{-m}
=\Bigl(\frac32\Bigr)^{r}-1 .
\]
For the first term, $n_k<2X$ and $2^{B_k}>2^{cr}\ge X$ by \eqref{eq:rlower},
so with \eqref{eq:rupper},
\[
\frac{3^rn_k}{2^{B_k}}<\frac{2X\cdot3^r}{X}=2\cdot3^r\le6\,X^{1-\eta}.
\]
Also $(3/2)^r\le\frac32X^{\kappa\log(3/2)}=\frac32X^{0.3562946\ldots}
\le\frac32X^{1-\eta}$.  Altogether
\[
1\le n_{k+r}<\tfrac{15}2X^{1-\eta}\qquad(k\in N_{\mathrm{high}}).
\]
By aperiodicity $k\mapsto n_{k+r}$ is injective, so $N_{\mathrm{high}}$
injects into the positive integers below $\tfrac{15}2X^{1-\eta}$, whence
$\#N_{\mathrm{high}}\le\tfrac{15}2X^{1-\eta}$.

\smallskip
Adding the two parts gives $V(X)\le11X^{1-\eta}$.
\end{proof}

\begin{remark}[why $c_*$ and why $\eta$ cannot be improved by this method]
\label{rem:opt}
Run the argument with free parameters $c\in(\alpha,2)$ and $\kappa>0$.  The
four contributions have exponents
\[
\underbrace{1-\kappa I(c)}_{\text{residue classes}},\quad
\underbrace{\kappa g(c)}_{\text{words}},\quad
\underbrace{1-\kappa\bigl(c\log2-\log3\bigr)}_{\text{growth}},\quad
\underbrace{\kappa\log(3/2)}_{\text{tail}} ,
\]
and $V(X)\ll X^{E}$ with $E$ the maximum of the four.  The fourth is never
binding.  Balancing the second against the smaller of the first and third
gives, after minimizing in $\kappa$,
\[
E(c)=\frac{g(c)}{g(c)+\min\bigl(I(c),\,c\log2-\log3\bigr)}
=\begin{cases}
\dfrac{g(c)}{c\log2-\log3+g(c)}, & g(c)\le\log3,\\[2ex]
\dfrac{g(c)}{c\log2}, & g(c)\ge\log3 .
\end{cases}
\]
On $(\alpha,c_*]$ the first branch decreases from $1$; on $[c_*,2)$ the second
branch increases, since $\frac{d}{dc}(g/c)=(cg'(c)-g(c))/c^2>0$ there.  The
minimum is therefore attained exactly where the two branches meet, that is at
$g(c)=\log3$, giving $c=c_*$, $\kappa=\kappa_*$ and $E=\alpha/c_*=1-\eta$.
A numerical two-parameter search over the pair $(c,\kappa)$ confirms this
optimum; see Appendix~\ref{app:numeric}.  Improving $\eta$ therefore requires
a different input, not a better choice of parameters.
\end{remark}

% ===========================================================================
\section{Reciprocal summability}
\label{sec:recip}

\begin{proof}[Proof of Corollary~\textup{\ref{cor:B}}]
Every orbit value lies in exactly one dyadic block $[2^q,2^{q+1})$ with
$q\ge0$, and each value there contributes at most $2^{-q}$ to the reciprocal
sum.  By Theorem~\ref{thm:A},
\[
\sum_{k\ge0}\frac1{n_k}
\le\sum_{q\ge0}V_0(2^q)\,2^{-q}
\le11\sum_{q\ge0}2^{q(1-\eta)}2^{-q}
=\frac{11}{1-2^{-\eta}}<464 .
\]
Since the series converges, $1/n_k\to0$, i.e.\ $n_k\to\infty$.  From
$2^{a_k}n_{k+1}=3n_k+1$ we get
$Y_{k+1}/Y_k=1+\frac1{3n_k}$ for $Y_k=2^{A_k}n_k/3^k$, so $(Y_k)$ is
increasing and
\[
Y_k=n_0\prod_{j<k}\Bigl(1+\frac1{3n_j}\Bigr)
\le n_0\exp\Bigl(\frac13\sum_{j\ge0}\frac1{n_j}\Bigr)<\infty ,
\]
whence $Y_k\uparrow L<\infty$ and $n_k=Y_k3^k2^{-A_k}\le L\,2^{\alpha k-A_k}$.
For the last statement, if $n_s=\min_{k\ge s}n_k$ then every $n_k$ with
$k\ge s$ lies in some $[2^qn_s,2^{q+1}n_s)$, $q\ge0$, and the same computation
with Theorem~\ref{thm:A} applied at $X=2^qn_s$ gives the bound
$11(1-2^{-\eta})^{-1}n_s^{-\eta}$.
\end{proof}

\begin{remark}
Corollary~\ref{cor:B} is stronger than the assertion that a divergent orbit
tends to infinity, and it is not implied by it: an orbit could tend to
infinity while returning to moderate values often enough to make
$\sum1/n_k$ diverge.  Theorem~\ref{thm:A} excludes that behaviour
quantitatively.  Note also that divergence is a \emph{conclusion} here, not a
hypothesis: aperiodicity alone forces $n_k\to\infty$.
\end{remark}

% ===========================================================================
\section{Repetitions and factor complexity}
\label{sec:symbolic}

The following proposition is the common engine of both parts of
Theorem~\ref{thm:CD}.  It converts any recurrence in the valuation word into
a height inequality.

\begin{proposition}[height bound for recurrences]
\label{prop:height}
Let $(n_k)$ be aperiodic with constant $L$ as in Corollary~\textup{\ref{cor:B}}.
Suppose $a_{u+i}=a_{u+p+i}$ for $0\le i<m$, where $u\ge0$, $p\ge1$, $m\ge1$,
and let $W=A_{u+m}-A_u$ be the weight of the repeated factor.  Then
\begin{equation}\label{eq:height}
W+1\ \le\ \log_2L+\log_2(3/2)\,(u+p).
\end{equation}
\end{proposition}

\begin{proof}
By Lemma~\ref{lem:collide}, $2^{W+1}\le|n_u-n_{u+p}|\le\max(n_u,n_{u+p})$,
the last step because both values are positive.  By
Corollary~\ref{cor:B} and $A_k\ge k$,
\[
n_k\le L\,2^{\alpha k-A_k}\le L\,2^{(\alpha-1)k}
\qquad(k\ge0),
\]
and $\alpha-1=\log_2(3/2)$.  Since the larger of the two indices is $u+p$,
\[
2^{W+1}\le L\,2^{\log_2(3/2)(u+p)},
\]
which is \eqref{eq:height}.
\end{proof}

\begin{proof}[Proof of Theorem~\textup{\ref{thm:CD}}]
\emph{The repetition exponent.}  Let $P=UV^tV'$ be a prefix of $a$, with
$u=|U|$, $p=|V|\ge1$, $t\ge1$ and $V'$ a prefix of $V$ of length $e$.  Then
$a_{u+i}=a_{u+p+i}$ for $0\le i<(t-1)p+e$, so Proposition~\ref{prop:height}
applies with $m=(t-1)p+e$ and gives
\[
A_{u+m}-A_u\le\log_2L+\log_2(3/2)(u+p).
\]
Every letter is at least $1$, so $A_{u+m}-A_u\ge m=(t-1)p+e$, and therefore
\[
|P|=u+tp+e=(u+p)+\bigl[(t-1)p+e\bigr]
\le(u+p)\bigl(1+\log_2(3/2)\bigr)+\log_2L
=\alpha\,(u+p)+\log_2L .
\]
Dividing by $u+p$ and letting $u+p\to\infty$ gives
$\operatorname{Dio}(a)\le\alpha=\log_23$.  (For prefixes with $u+p$ bounded
the ratio is bounded anyway, so the supremum defining
$\operatorname{Dio}$ is governed by the limit.)

\emph{The factor complexity.}  Fix $m\ge1$ and let $J\ge1$.  If
$p_a(m)<J$ then two of the $J$ factors of length $m$ occurring at positions
$0,1,\dots,J-1$ coincide, say at $i<j<J$.  Proposition~\ref{prop:height}
with $u=i$, $p=j-i$ gives
\[
m\le A_{i+m}-A_i\le\log_2L+\log_2(3/2)\,j\le\log_2L+\log_2(3/2)\,J .
\]
Contrapositively, if $J<(m-\log_2L)/\log_2(3/2)$ then $p_a(m)\ge J$.  Taking
$J=\lceil(m-\log_2L)/\log_2(3/2)\rceil-1$ yields
\[
p_a(m)\ge\frac{m-\log_2L}{\log_2(3/2)}-1 .
\]
Dividing by $m$ and letting $m\to\infty$ gives the stated $\liminf$.  A
Sturmian word has $p(m)=m+1$ and a quasi-Sturmian word has $p(m)=m+k$ for
all large $m$; since $1/\log_2(3/2)>1$, neither is compatible with the
displayed bound for large $m$.
\end{proof}

\begin{remark}
The bound $p_a(m)\ge1.7095\ldots m-O(1)$ should be compared with the
Morse--Hedlund theorem \cite{morsehedlund1940}, which only excludes
$p_a(m)\le m$ for an aperiodic word.  The gain is the factor
$1/\log_2(3/2)$, and it comes entirely from the arithmetic of the map rather
than from combinatorics on words.  The constant is essentially the same
$\rho=\log2/\log(3/2)$ that governs the recharge clock of the orbit.
\end{remark}

\begin{remark}[a stronger form that we do \emph{not} obtain]
It is tempting to replace $m$ in Theorem~\ref{thm:CD} by $A_m=m+\sum_{k<m}(a_k-1)$,
i.e.\ by the weight of the \emph{initial} factor of length $m$.  The
pigeonhole argument gives the weight $A_{i+m}-A_i$ of the recurring factor,
which is only bounded below by $m$; there is no reason for an arbitrary
window to be as heavy as the initial one.  The stronger form is valid when
the recurrence involves position $i=0$, but not in general.
\end{remark}

% ===========================================================================
\section{What these bounds do not give}
\label{sec:limits}

We list the limitations explicitly, since the results above are easy to
overread.

\begin{enumerate}
  \item \emph{No contradiction.}  All four statements are consistent with the
        existence of a divergent orbit.  They constrain such an orbit; they do
        not exclude one.  Theorem~\ref{thm:A} allows $X^{0.9653\ldots}$ visits
        to a shell, which is very far from the $O(1)$ that a
        counting contradiction would require.

  \item \emph{No control of time.}  $V_s(X)$ counts how \emph{often} an orbit
        is small, not how \emph{late}.  Nothing above bounds last-exit times,
        or the rate at which $\sum_{j\ge k}1/n_j$ tends to $0$ as a function
        of $k$ rather than of $n_s$.

  \item \emph{The exponent is method-limited.}  By Remark~\ref{rem:opt} the
        value $\eta=0.0346\ldots$ is the exact optimum of this argument, so
        improvements must change the input, for instance by exploiting
        $3$-adic integrality of inverse branches, which is used nowhere here.

  \item \emph{Complexity is a weak filter.}  A bound of the form
        $p_a(m)\gtrsim1.71m$ excludes only the lowest-complexity words.
        Words of exponential factor complexity and repetition exponent $1$
        satisfy every constraint proved here, and such words are exactly what
        a divergent orbit is expected to produce.

  \item \emph{Cycles are untouched.}  Every result assumes aperiodicity.
\end{enumerate}

\begin{remark}[relation to prior work]
The ingredients are individually classical: the cylinder structure of
Lemma~\ref{lem:cyl} is the parity-vector correspondence of Terras
\cite{terras1976} and Bernstein--Lagarias \cite{bernsteinlagarias1996}, and
the exponential sum of valuations is the object estimated in
\cite{tao2022}.  What appears not to be standard is the combination:
applying these to a \emph{single} orbit, using non-collision of orbit values
in place of a density or measure hypothesis.  We have not located the
statements of Theorem~\ref{thm:A}, Corollary~\ref{cor:B} or
Theorem~\ref{thm:CD} in the literature, but the $3x+1$ bibliography is large
\cite{lagarias1985} and our search was not exhaustive; readers aware of a
prior occurrence are asked to inform us.
\end{remark}

% ===========================================================================
\appendix
\section{Numerical verification}
\label{app:numeric}

Every lemma was checked numerically before the proofs were written.  The
checks are finite and prove nothing; they were used to catch errors.

\begin{itemize}
  \item \textbf{Lemma~\ref{lem:cyl}.}  For all odd $n<2^{14}$ and all
        $r\le8$, the $\num{20964}$ realized words were grouped by value; every
        word occupied exactly one class modulo $2^{B+1}$.  No violation.
  \item \textbf{Lemma~\ref{lem:collide}.}  Over $400$ random starts below
        $10^{12}$ and factors of length $6$, all $\num{3503}$ coincident pairs
        satisfied $n_i\equiv n_j\pmod{2^{W+1}}$.  No violation.
  \item \textbf{Endpoint bound (proof of Theorem~\ref{thm:A}).}  For
        $\num{3000}$ random blocks the inequality
        $n_{k+r}\le2X(3/2^{c})^{r}+(3/2)^{r}$ held throughout.
  \item \textbf{Optimality (Remark~\ref{rem:opt}).}  A grid search over
        $(c,\kappa)\in(\alpha,2)\times(0,2]$ with all four exponents returned
        the minimum $0.9654$ at $c=1.6418$, $\kappa=0.8780$, matching
        $c_*,\kappa_*$ to grid resolution.
  \item \textbf{Proposition~\ref{prop:height}.}  Over $300$ random starts
        below $10^{14}$, all $\num{640362}$ recurrences in the valuation words
        satisfied \eqref{eq:height} with $L$ computed exactly as
        $\max_kY_k$ in rational arithmetic.  No violation; minimum slack
        $33.1$.
  \item \textbf{Theorem~\ref{thm:CD}.}  In $\num{1716}$ cases the complexity
        bound held, with observed $p(m)/m$ between $6.75$ and $16$.
\end{itemize}

Since a genuinely aperiodic orbit is not available for testing, the checks of
Proposition~\ref{prop:height} and Theorem~\ref{thm:CD} were run on the finite
transients of convergent orbits, with $L$ replaced by $\max_kY_k$ over the
transient; this is the correct finite analogue, because the proofs use only
$n_k\le L2^{\alpha k-A_k}$ on the range of indices involved.

\printbibliography[title={References}]

\end{document}
